% Figure: GS QKD 018 reference scenario (adapted from GS QKD 018, Fig. 1)
\begin{figure}[t]
\centering
\begin{tikzpicture}
  \node[box,fill=blue!8,minimum width=3.3cm] (orch) {SDN orchestrator};
  \node[box,fill=green!8,minimum width=3.3cm,below left=12mm and 6mm of orch] (cq)
       {SDN controller\\QKD network};
  \node[box,fill=orange!10,minimum width=3.3cm,below right=12mm and 6mm of orch] (co)
       {SDN controller\\OTN};
  \node[box,fill=green!4,below=10mm of cq,minimum width=4.6cm] (nq)
       {SD-QKD nodes (agent, KME)};
  \node[box,fill=orange!5,below=10mm of co,minimum width=4.6cm] (no)
       {OTN nodes (SAE)};
  \draw[<->,very thick,red!70!black] (orch) -- node[lbl]{GS QKD 018} (cq);
  \draw[<->,thick] (orch) -- node[lbl]{OTN NBI (out of scope)} (co);
  \draw[<->,thick] (cq) -- node[lbl]{GS QKD 015} (nq);
  \draw[<->,thick] (co) -- (no);
  \draw[<->,dashed] (nq) -- node[lbl]{QKD 004 / 014} (no);
\end{tikzpicture}
\caption{GS~QKD~018 reference scenario (adapted
from~\cite[Fig.~1]{qkd018}). The standardised interface is in red; the
dashed line is key delivery from KMEs to SAEs, which happens inside the
same security perimeter.}
\label{fig:scenario}
\end{figure}
